\chapter{Discussion}\label{ch:discussion}

The experiments of Chapter~\ref{ch:experiments} support the use of
$k$-means initialisation, and they raise several questions that this chapter
discusses. The estimator is that of Chapter~\ref{ch:estimation}, computed by
the algorithm of Chapter~\ref{ch:algorithms}.

\section{Speed and accuracy}\label{sec:discussion-speed}

Table~\ref{tab:recovery} shows that accuracy improves with sample size under
every strategy. The cost per iteration scales with $nkd^2$, which is dominated
by the covariance updates of Equation~\eqref{eq:upd-sigma}. Parameter counts from
Equation~\eqref{eq:count} give the complexity of the model itself.

\lipsum[41-42]

\section{Limitations}\label{sec:limitations}

The analysis assumes Gaussian components, and real data may be heavy-tailed.
Identifiability, discussed in Section~\ref{sec:identifiability}, also depends on
separation between components, which the synthetic designs control only
partially.

\lipsum[43-44]

\section{Future work}\label{sec:future}

Extensions to non-Gaussian components and to streaming updates are natural next
steps. Chapter~\ref{ch:intro} lists the questions that motivated this work.

\lipsum[45-46]

\section{Further remarks}\label{sec:further-discussion}

The remarks in this section restate the main point of the chapter in more
general terms, so that the argument can be read without the surrounding
derivations.

\lipsum[123-134]
